% source: manti.01.ma, 1ère SM, devoir surveillé n°1 (Drive 0B6pfSJcEVng0SWFYZWVXbDFtZHM). Ex.3 logique/arithmétique (hors thème)
\ex{Fonctions}{SWFY}{Ex. 1}{6}
On considère la fonction $f$ définie par $f(x)=\dfrac{x^3}{x^2-3}$.
\begin{enumerate}
\item Déterminer l'ensemble de définition de $f$.
\item Développer le produit $(x-3)^2(2x+3)$ et montrer que $f$ admet un extremum en $x_0=3$ en précisant sa nature.
\end{enumerate}
\fin
\ex{Fonctions}{SWFY}{Ex. 2}{10}
Soit $f$ la fonction définie par $f(x)=\dfrac{x^2-3}{x-2}$.
\begin{enumerate}
\item \begin{enumerate}
\item Montrer que $\dfrac{f(x)-f(y)}{x-y}=\dfrac{xy-2x-2y+3}{(x-2)(y-2)}$.
\item Montrer que $f$ est croissante sur l'intervalle $]-\infty,1]$, puis étudier la monotonie de $f$ sur l'intervalle $[1,2[$.
\end{enumerate}
\item On pose $g(x)=\sqrt{x+4}$ et $h(x)=\dfrac{x+1}{x}\left(\sqrt{x+4}+2\right)$.
\begin{enumerate}
\item Déterminer $D_h$ et $D_{f\circ g}$.
\item Vérifier que $\big(\forall x\in[-4,-3]\big)\ h(x)=(f\circ g)(x)$.
\item Étudier la monotonie de la fonction $h$ sur l'intervalle $[-4,-3]$.
\end{enumerate}
\end{enumerate}
\fin
